% Generated draft: requires mathematical and source review.
\documentclass[11pt]{article}
\usepackage[margin=25mm]{geometry}
\usepackage{fontspec}
\setmainfont{Latin Modern Roman}
\usepackage{amsmath,amssymb,mathtools,unicode-math}
\setmathfont{Latin Modern Math}
\usepackage{enumitem,needspace,xurl,hyperref}
\setlength{\parindent}{0pt}
\setlength{\parskip}{5pt}
\setlength{\emergencystretch}{3em}
\newcommand{\N}{\mathbb{N}}
\newcommand{\Z}{\mathbb{Z}}
\newcommand{\Q}{\mathbb{Q}}
\newcommand{\R}{\mathbb{R}}
\newcommand{\C}{\mathbb{C}}
\newcommand{\F}{\mathbb{F}}
\newcommand{\A}{\mathbb{A}}
\newcommand{\PP}{\mathbb{P}}
\newcommand{\E}{\mathbb{E}}
\newcommand{\eps}{\varepsilon}
\newcommand{\dd}{\,\mathrm{d}}
\begin{document}
\section*{UnsolvedMath ID: 30005239}
\textbf{Source problem code:} catalogue entry
\subsection{Definitions and mathematical statement}
The supplied fragment concerns a differential operator \(P(D)\) with constant coefficients and polynomial symbol, transcribed from the damaged summation as
\[
 p(\xi)=\sum_{|\alpha|\leq N}c_\alpha\xi^\alpha.
\]
In this notation, \(\alpha=(\alpha_1,\ldots,\alpha_n)\) is a multi-index of nonnegative integers, \(|\alpha|=\alpha_1+\cdots+\alpha_n\), and \(\xi^\alpha=\prod_{j=1}^n\xi_j^{\alpha_j}\). The dimension \(n\), order bound \(N\), coefficient field, underlying domain, and convention for \(D\) are not specified in the fragment. The displayed polynomial is a transcription of the constant-coefficient example, not a supplied formulation of the variable-coefficient problem. The equations mentioned are
\[
 P(D)u=0
 \qquad\text{or}\qquad
 P(D)u=\text{``small''}.
\]
No norm, topology, scale, or parameter defining ``small'' is supplied.

\emph{Recoverable question.} For partial differential equations with nonconstant coefficients, can the phase-space maps referred to in the surrounding discussion be quantized so that they act at the level of operators?

Here ``quantized'' retains the source's informal meaning of passing from phase-space maps to operators; the fragment supplies no mathematical definition or required operator identity. It does not identify the maps, their domains, the class of differential operators, or the function spaces on which a quantization should act. In particular, ``any of the phase space maps'' does not determine a universal quantifier over all phase-space maps. Whether the question concerns local or global constructions, exact identities or asymptotic relations, or any prescribed boundary conditions is unresolved. Consequently, the supplied evidence does not determine a precise existence theorem or conjecture.

The supplied status is \texttt{partially\_solved}, and the statement is marked \texttt{unrecoverable}. The supplied dated literature assessment says that semiclassical Fourier integral operators give local quantizations when the characteristic set is a graph, while the fragment does not establish a construction for arbitrary variable-coefficient equations or global gluing. A separate supplied verification note says that the official report immediately answers the question using a semiclassical Fourier-integral-operator construction and that it is therefore not an open problem. These statements are supplied curation, not independently verified mathematical conclusions; their scopes cannot be reconciled from the fragment alone.
\subsection{Short English statement}
The fragment asks whether phase-space maps for variable-coefficient partial differential equations can be represented by operators. It omits the maps, hypotheses, and meaning of quantization needed for a precise mathematical statement. The supplied partial-solution assessment and the separate claim that the report already answers the question require reconciliation.
\subsection{Sources}
\begin{itemize}
\item Local UnsolvedMath registry: supplied complete source record, permanent ID \(30005239\), title \emph{Phase-Space Quantization for Variable-Coefficient PDEs}, exact source problem code \texttt{catalogue entry}; collection: Oberwolfach Reports Open Problems. The record gives proposed year \(2022\) and names no proposer.
\item Supplied report citation: \emph{At the Interface between Semiclassical Analysis and Numerical Analysis of Wave Scattering Problems} (2023), Oberwolfach Reports, pp.~2511--2587, DOI \texttt{10.4171/owr/2022/43}. Supplied source URL: \url{https://doi.org/10.4171/owr/2022/43}. The background identifies Oberwolfach Reports as records of mathematical workshops and their open-problem discussions; no passage-level location for this fragment is supplied.
\item Supplied literature source: \emph{Oberwolfach Report 43/2022}, \url{https://doi.org/10.4171/owr/2022/43}. Literature checked at: \texttt{2026-08-22}. This date and the associated partial-solution assessment are supplied dated curation, not independently verified. The background explicitly calls the assessment dated literature triage derived from supplied OWR curation and requests independent verification of time-sensitive status and cited claims.
\item The supplied statement-verification field claims that the report immediately answers the question through a semiclassical Fourier-integral-operator construction. This attribution is preserved as a supplied claim; the report has not been independently consulted for this draft.
\end{itemize}
\end{document}
